% !TEX program = xelatex
% Arabic document with fully vocalized text (tashkeel / harakat).
% Shows fatha, damma, kasra, sukun, shadda, tanwin in context.
\documentclass[12pt, a4paper]{article}

\usepackage{fontspec}
\usepackage{polyglossia}
\setmainlanguage[locale=mashriq]{arabic}
\setotherlanguage{english}

\newfontfamily\arabicfont[Script=Arabic]{Amiri}

\begin{document}

\title{التشكيلُ والحركاتُ في النصِّ العربيّ}
\author{د. نبراس أبو الذهب}
\date{2025}
\maketitle

\section{الحركاتُ الأساسية}

يَستعرضُ هذا المثالُ الحركاتِ كاملةً في نصٍّ عربيٍّ مُشكَّلٍ بالكامل.

\subsection{الفتحةُ والضمةُ والكسرة}

بَرَعَ الْعَرَبُ في فَصاحَةِ اللِّسانِ، وَأَحْسَنُوا التَّعْبِيرَ عَنِ
الْمَعانِي الدَّقِيقَةِ. كَتَبَ الْعُلَماءُ في الْفِقْهِ وَالنَّحْوِ
وَالْبَلاغَةِ، وَتَفَرَّدُوا بِعِلْمِ التَّشْكِيلِ الَّذِي يَحْفَظُ
الْقِراءَةَ مِنَ اللَّحْنِ.

\subsection{السكونُ والشدّةُ والتنوين}

إِنَّ التَّشْكِيلَ عِلْمٌ ثابِتٌ، يَضْبِطُ النُّطْقَ وَيُوَضِّحُ
الْمَعْنَى. اِقْتَضَتْ قَواعِدُهُ رَسْمَ الْحَرَكاتِ فَوْقَ الْحُروفِ
وَتَحْتَها، وَرَسْمَ الشَّدَّةِ وَالسُّكُونِ وَالتَّنْوِينِ.
الْعِلْمُ نورٌ، وَالنُّورُ هِدايَةٌ، وَالْهِدايَةُ سَعادَةٌ.

\subsection{نصٌّ مُشكَّلٌ بالكامل}

قَالَ تَعالَى: \LR{﴿}اقْرَأْ بِاسْمِ رَبِّكَ الَّذِي خَلَقَ\LR{﴾}
خَلَقَ الْإِنْسانَ مِنْ عَلَقٍ، اقْرَأْ وَرَبُّكَ الْأَكْرَمُ،
الَّذِي عَلَّمَ بِالْقَلَمِ، عَلَّمَ الْإِنْسانَ ما لَمْ يَعْلَمْ.

\end{document}
